\documentclass[12pt,a4paper]{article}
\usepackage[utf8]{inputenc}
\usepackage[arabic,english]{babel}
\usepackage{graphicx}
\usepackage{hyperref}
\usepackage{geometry}
\usepackage{listings}
\usepackage{color}

\geometry{margin=1in}

\title{UMBB MobGuide: Mobile Project Report}
\author{Group Members: Amine \& Team}
\date{April 2026}

\begin{document}

\maketitle

\section{Introduction}
The \textbf{UMBB MobGuide} is a digital guide developed for the University of M'hamed Bougara of Boumerdès (UMBB). The application aims to provide students, researchers, and visitors with a centralized platform to explore the various faculties and departments of the university, access detailed information, and facilitate direct contact and navigation.

\section{Application Components}
The application is built using the Android SDK with Java and XML, following a Model-View-Controller (MVC) architecture.

\subsection{Activities and Interfaces}
\begin{itemize}
    \item \textbf{MainActivity}: Serves as the landing page, displaying a general presentation of the university and its history.
    \item \textbf{FacultyListActivity}: Displays a comprehensive list of all faculties using a ListView. It includes a SearchView for filtering.
    \item \textbf{DepartmentListActivity}: Shows the departments associated with the selected faculty.
    \item \textbf{DetailActivity}: Provides in-depth information about a specific faculty or department, including specialties and contact actions.
\end{itemize}

\subsection{Data Management}
\begin{itemize}
    \item \textbf{DataManager}: A utility class that manages data using dynamic \texttt{ArrayLists}, simulating a database environment.
    \item \textbf{Models}: \texttt{Faculty} and \texttt{Department} classes define the structure of university entities.
\end{itemize}

\section{Features Implemented}
The following features have been successfully integrated:
\begin{enumerate}
    \item \textbf{University Presentation}: Detailed textual description of UMBB.
    \item \textbf{Faculty \& Department Browsing}: Hierarchical navigation through university structures.
    \item \textbf{Search Functionality}: Real-time search integrated into the list activities.
    \item \textbf{Multilingual Interface}: Support for both English and Arabic.
    \item \textbf{Interactive Actions}:
    \begin{itemize}
        \item \textbf{Call}: Direct dialing to the administration.
        \item \textbf{SMS}: Pre-filled recipient for quick inquiries.
        \item \textbf{Email}: Automated composition of inquiry emails.
        \item \textbf{Google Maps}: GPS navigation to the physical location.
    \end{itemize}
\end{enumerate}

\section{Task Distribution}
The project was completed in a collaborative environment with the following distribution of roles:
\begin{itemize}
    \item \textbf{Lead Developer}: Implementation of core logic, activities, and intent handling.
    \item \textbf{UI/UX Designer}: Creation of XML layouts, custom font integration, and SVG icon design.
    \item \textbf{Data Analyst}: Information gathering about UMBB and data modeling.
    \item \textbf{Technical Writer}: Production of this LaTeX documentation and the presentation walkthrough.
\end{itemize}

\section{Conclusion}
The UMBB MobGuide successfully fulfills the pedagogical objectives of the mobile development module. It demonstrates practical knowledge of Android components, lifecycle management, and user-centric design.

\end{document}
